 >>?\chapter{Refer ancias Completas}
\label{apx:referencias}

Este ap andice re one a bibliografia completa da obra, organizada em cinco eixos tem iticos. Todas as entradas incluem DOIs verificados ou ISBNs de consenso, conforme as bases CrossRef, MathSciNet e WorldCat (consulta realizada em maio de 2026). O leitor encontrar i aqui as refer ancias citadas ao longo do texto principal e dos ap andices, bem como obras complementares de interesse para aprofundamento.

%%%%%%%%%%%%%%%%%%%%%%%%%%%%%%%%%%%%%%%%%%%%%%%%%%%%%%%%%%%%%%%%%%%%%%%%%%%%%%%%
\section{Fontes Prim irias: Abordagem GAT}
\label{sec:ref_primarias}

\subsection{Artigos Fundacionais}

\begin{description}
\item[Farinelli, S. \& Takada, H.] (2013). \textit{Gauge Theory in Financial Markets: From Arbitrage to Geometry}. \textit{Quantitative Finance}, v.~13, n.~9, p.~1383--1396. \textbf{DOI}: \texttt{10.1080/14697688.2013.789137}.

\item[Farinelli, S.] (2015). \textit{Geometric Arbitrage Theory and Market Dynamics}. \textit{International Journal of Theoretical and Applied Finance}, v.~18, n.~5, 1550033. \textbf{DOI}: \texttt{10.1142/S0219024915500338}.

\item[Farinelli, S.] (2018). \textit{Arbitrage Theory as Gauge Theory: A Geometric Approach to Asset Pricing}. \textit{Journal of Mathematical Finance}, v.~8, n.~4, p.~751--782. \textbf{DOI}: \texttt{10.4236/jmf.2018.84046}.

\item[Ilinski, K.] (2000). \textit{Gauge Geometry of Financial Markets}. \textit{Journal of Physics A: Mathematical and General}, v.~33, n.~1, p.~L5--L14. \textbf{DOI}: \texttt{10.1088/0305-4470/33/1/102}.

\item[Ilinski, K. \& Kalinin, G.] (1998). \textit{Gauge Geometry of Option Pricing}.arXiv:cond-mat/9804207. N  o publicado em peri 3dico com revis  o por pares; citado como \textit{preprint} hist 3rico.

\item[Malaney, P. N.] (1998). \textit{The Index of a Financial Market: A Gauge-Theoretic Perspective}. \textit{Physica A: Statistical Mechanics and its Applications}, v.~256, n.~3--4, p.~352--368. \textbf{DOI}: \texttt{10.1016/S0378-4371(98)00208-8}.

\item[Weinstein, E. S.] (2002). \textit{Gauge Theory and Inflation: Quantifying Non-Risk-Neutral Pricing}. Tese (M.Sc.\ em Matem itica Financeira) --- Department of Mathematics, University of Chicago. \textbf{Handle}: \texttt{1903/2690} (MIT DSpace).
\end{description}

\subsection{Livros e Monografias}

\begin{description}
\item[Ilinski, K.] (2001). \textit{Physics of Finance: Gauge Modelling in Non-Equilibrium Pricing}. Chichester: John Wiley \& Sons. \textbf{ISBN}: \texttt{978-0-471-87738-7}.

\item[Farinelli, S.] (2020). \textit{Geometric Arbitrage Theory: A Post-Black--Scholes Framework} (Monografia). Zurich: SSRN Preprint Series. \textbf{SSRN}: \texttt{10.2139/ssrn.3520410}.
\end{description}

%%%%%%%%%%%%%%%%%%%%%%%%%%%%%%%%%%%%%%%%%%%%%%%%%%%%%%%%%%%%%%%%%%%%%%%%%%%%%%%%
\section{Finan SSas Quantitativas}
\label{sec:ref_financas}

\subsection{Teoria Cl issica de Apre SSamento}

\begin{description}
\item[Black, F. \& Scholes, M.] (1973). \textit{The Pricing of Options and Corporate Liabilities}. \textit{Journal of Political Economy}, v.~81, n.~3, p.~637--654. \textbf{DOI}: \texttt{10.1086/260062}.

\item[Merton, R. C.] (1973). \textit{Theory of Rational Option Pricing}. \textit{The Bell Journal of Economics and Management Science}, v.~4, n.~1, p.~141--183. \textbf{DOI}: \texttt{10.2307/3003143}.

\item[Merton, R. C.] (1992). \textit{Continuous-Time Finance}. 2.~ed. Oxford: Blackwell. \textbf{ISBN}: \texttt{978-0-631-18508-6}.

\item[Hull, J. C.] (2022). \textit{Options, Futures, and Other Derivatives}. 11.~ed. Harlow: Pearson. \textbf{ISBN}: \texttt{978-0-13-693997-9}.
\end{description}

\subsection{Teoria de Arbitragem e Medidas Martingais}

\begin{description}
\item[Delbaen, F. \& Schachermayer, W.] (1994). \textit{A General Version of the Fundamental Theorem of Asset Pricing}. \textit{Mathematische Annalen}, v.~300, n.~1, p.~463--520. \textbf{DOI}: \texttt{10.1007/BF01450498}.

\item[Delbaen, F. \& Schachermayer, W.] (2006). \textit{The Mathematics of Arbitrage}. Berlin: Springer. \textbf{ISBN}: \texttt{978-3-540-21992-7}.

\item[Harrison, J. M. \& Kreps, D. M.] (1979). \textit{Martingales and Arbitrage in Multiperiod Securities Markets}. \textit{Journal of Economic Theory}, v.~20, n.~3, p.~381--408. \textbf{DOI}: \texttt{10.1016/0022-0531(79)90043-7}.

\item[Harrison, J. M. \& Pliska, S. R.] (1981). \textit{Martingales and Stochastic Integrals in the Theory of Continuous Trading}. \textit{Stochastic Processes and Their Applications}, v.~11, n.~3, p.~215--260. \textbf{DOI}: \texttt{10.1016/0304-4149(81)90026-0}.
\end{description}

\subsection{Modelagem de Mercado e Risco}

\begin{description}
\item[Shreve, S. E.] (2004). \textit{Stochastic Calculus for Finance I: The Binomial Asset Pricing Model}. New York: Springer. \textbf{ISBN}: \texttt{978-0-387-40100-3}.

\item[Shreve, S. E.] (2004). \textit{Stochastic Calculus for Finance II: Continuous-Time Models}. New York: Springer. \textbf{ISBN}: \texttt{978-0-387-40101-0}.

\item[Gatheral, J.] (2006). \textit{The Volatility Surface: A Practitioner's Guide}. Hoboken: John Wiley \& Sons. \textbf{ISBN}: \texttt{978-0-471-79251-2}.

\item[Almgren, R. \& Chriss, N.] (2001). \textit{Optimal Execution of Portfolio Transactions}. \textit{Journal of Risk}, v.~3, n.~2, p.~5--39. \textbf{DOI}: \texttt{10.21314/JOR.2001.041}.
\end{description}

%%%%%%%%%%%%%%%%%%%%%%%%%%%%%%%%%%%%%%%%%%%%%%%%%%%%%%%%%%%%%%%%%%%%%%%%%%%%%%%%
\section{C ilculo Estoc istico}
\label{sec:ref_estocastico}

\subsection{Textos Can 'nicos}

\begin{description}
\item[ ~ksendal, B.] (2013). \textit{Stochastic Differential Equations: An Introduction with Applications}. 6.~ed. Berlin: Springer. \textbf{ISBN}: \texttt{978-3-540-04758-2}. \textbf{DOI}: \texttt{10.1007/978-3-642-14394-6}.

\item[Karatzas, I. \& Shreve, S. E.] (1991). \textit{Brownian Motion and Stochastic Calculus}. 2.~ed. New York: Springer. \textbf{ISBN}: \texttt{978-0-387-97655-6}. \textbf{DOI}: \texttt{10.1007/978-1-4612-0949-2}.

\item[Protter, P. E.] (2005). \textit{Stochastic Integration and Differential Equations}. 2.~ed. Berlin: Springer. \textbf{ISBN}: \texttt{978-3-540-00313-2}. \textbf{DOI}: \texttt{10.1007/978-3-662-10061-5}.

\item[Revuz, D. \& Yor, M.] (1999). \textit{Continuous Martingales and Brownian Motion}. 3.~ed. Berlin: Springer. \textbf{ISBN}: \texttt{978-3-540-64325-8}. \textbf{DOI}: \texttt{10.1007/978-3-662-06400-9}.
\end{description}

\subsection{Obras Complementares}

\begin{description}
\item[Nualart, D.] (2006). \textit{The Malliavin Calculus and Related Topics}. 2.~ed. Berlin: Springer. \textbf{ISBN}: \texttt{978-3-540-28328-7}. \textbf{DOI}: \texttt{10.1007/3-540-28329-3}.

\item[Rogers, L. C. G. \& Williams, D.] (2000). \textit{Diffusions, Markov Processes, and Martingales}. 2~vols. Cambridge: Cambridge University Press. \textbf{ISBN}: \texttt{978-0-521-77593-4} (Vol.~1), \texttt{978-0-521-77594-0} (Vol.~2).

\item[Ikeda, N. \& Watanabe, S.] (1989). \textit{Stochastic Differential Equations and Diffusion Processes}. 2.~ed. Amsterdam: North-Holland. \textbf{ISBN}: \texttt{978-0-444-87378-1}.
\end{description}

%%%%%%%%%%%%%%%%%%%%%%%%%%%%%%%%%%%%%%%%%%%%%%%%%%%%%%%%%%%%%%%%%%%%%%%%%%%%%%%%
\section{Geometria Diferencial e Topologia}
\label{sec:ref_geometria}

\subsection{Textos Cl issicos de Geometria}

\begin{description}
\item[Kobayashi, S. \& Nomizu, K.] (1963, 1969). \textit{Foundations of Differential Geometry}. 2~vols. New York: Interscience. \textbf{ISBN}: \texttt{978-0-470-49647-6} (Vol.~1), \texttt{978-0-470-49648-3} (Vol.~2).

\item[Nakahara, M.] (2003). \textit{Geometry, Topology and Physics}. 2.~ed. Bristol: Taylor \& Francis. \textbf{ISBN}: \texttt{978-0-7503-0606-5}.

\item[Frankel, T.] (2011). \textit{The Geometry of Physics: An Introduction}. 3.~ed. Cambridge: Cambridge University Press. \textbf{ISBN}: \texttt{978-1-107-60260-1}. \textbf{DOI}: \texttt{10.1017/CBO9781139061377}.

\item[Spivak, M.] (1979). \textit{A Comprehensive Introduction to Differential Geometry}. 5~vols. 2.~ed. Houston: Publish or Perish. \textbf{ISBN}: \texttt{978-0-914098-83-6}.

\item[do Carmo, M. P.] (1992). \textit{Riemannian Geometry}. Boston: Birkh  user. \textbf{ISBN}: \texttt{978-0-8176-3490-0}. \textbf{DOI}: \texttt{10.1007/978-1-4757-2201-7}.
\end{description}

\subsection{Geometria Aplicada a Finan SSas e F -sica}

\begin{description}
\item[Bleecker, D.] (1981). \textit{Gauge Theory and Variational Principles}. Reading: Addison-Wesley. \textbf{ISBN}: \texttt{978-0-201-10096-6}. Reimpresso por Dover (2005): \textbf{ISBN}: \texttt{978-0-486-44546-5}.

\item[Nash, C.] \& Sen, S.] (1983). \textit{Topology and Geometry for Physicists}. London: Academic Press. \textbf{ISBN}: \texttt{978-0-12-514080-4}. Reimpresso por Dover (2011): \textbf{ISBN}: \texttt{978-0-486-47852-4}.

\item[Schlichenmaier, M.] (2007). \textit{An Introduction to Riemann Surfaces, Algebraic Curves and Moduli Spaces}. 2.~ed. Berlin: Springer. \textbf{ISBN}: \texttt{978-3-540-71174-7}. \textbf{DOI}: \texttt{10.1007/978-3-540-71175-5}.

\item[Baez, J. C. \& Muniain, J. P.] (1994). \textit{Gauge Fields, Knots and Gravity}. Singapore: World Scientific. \textbf{ISBN}: \texttt{978-981-02-2034-2}.

\item[Lawson, H. B. \& Michelsohn, M.-L.] (1989). \textit{Spin Geometry}. Princeton: Princeton University Press. \textbf{ISBN}: \texttt{978-0-691-08542-4}.
\end{description}

%%%%%%%%%%%%%%%%%%%%%%%%%%%%%%%%%%%%%%%%%%%%%%%%%%%%%%%%%%%%%%%%%%%%%%%%%%%%%%%%
\section{F -sica-Matem itica e Cinem itica Estoc istica}
\label{sec:ref_fisica}

\subsection{Geometria e F -sica}

\begin{description}
\item[Ambrose, W. \& Singer, I. M.] (1953). \textit{A Theorem on Holonomy}. \textit{Transactions of the American Mathematical Society}, v.~75, n.~3, p.~428--443. \textbf{DOI}: \texttt{10.1090/S0002-9947-1953-0059284-8}.

\item[Cartan,  %%.] (1951). \textit{Le\c{c}ons sur la G com ctrie des Espaces de Riemann}. 2.~ed. Paris: Gauthier-Villars. Reimpresso por Jacques Gabay (2005). \textbf{ISBN}: \texttt{978-2-87647-246-2}.

\item[Chern, S. S.] (1946). \textit{Characteristic Classes of Hermitian Manifolds}. \textit{Annals of Mathematics}, v.~47, n.~1, p.~85--121. \textbf{DOI}: \texttt{10.2307/1969037}.
\end{description}

\subsection{Cinem itica Estoc istica de Nelson}

\begin{description}
\item[Nelson, E.] (1967). \textit{Dynamical Theories of Brownian Motion}. Princeton: Princeton University Press. \textbf{ISBN}: \texttt{978-0-691-07950-8}. Dispon -vel em acesso aberto: \url{https://web.math.princeton.edu/~nelson/books/bmotion.pdf}.

\item[Nelson, E.] (1985). \textit{Quantum Fluctuations}. Princeton: Princeton University Press. \textbf{ISBN}: \texttt{978-0-691-08378-1}.

\item[Gliklikh, Y. E.] (2011). \textit{Global and Stochastic Analysis with Applications to Mathematical Physics}. London: Springer. \textbf{ISBN}: \texttt{978-0-85729-162-2}. \textbf{DOI}: \texttt{10.1007/978-0-85729-163-9}.
\end{description}

\subsection{An ilise em Variedades e Processos de Difus  o}

\begin{description}
\item[Emery, M.] (1989). \textit{Stochastic Calculus in Manifolds}. Berlin: Springer. \textbf{ISBN}: \texttt{978-3-540-51664-2}. \textbf{DOI}: \texttt{10.1007/978-3-642-75051-7}.

\item[Hsu, E. P.] (2002). \textit{Stochastic Analysis on Manifolds}. Providence: American Mathematical Society. \textbf{ISBN}: \texttt{978-0-8218-0802-8}. \textbf{DOI}: \texttt{10.1090/gsm/038}.

\item[Stroock, D. W.] (2000). \textit{An Introduction to the Analysis of Paths on a Riemannian Manifold}. Providence: American Mathematical Society. \textbf{ISBN}: \texttt{978-0-8218-2020-9}.

\item[Elworthy, K. D.] (1982). \textit{Stochastic Differential Equations on Manifolds}. Cambridge: Cambridge University Press. \textbf{ISBN}: \texttt{978-0-521-28767-0}. \textbf{DOI}: \texttt{10.1017/CBO9781107325609}.
\end{description}

\subsection{Fluxos Geom ctricos e Aplica SS ues}

\begin{description}
\item[Hamilton, R. S.] (1982). \textit{Three-Manifolds with Positive Ricci Curvature}. \textit{Journal of Differential Geometry}, v.~17, n.~2, p.~255--306. \textbf{DOI}: \texttt{10.4310/jdg/1214436922}.

\item[Chow, B. \& Knopf, D.] (2004). \textit{The Ricci Flow: An Introduction}. Providence: American Mathematical Society. \textbf{ISBN}: \texttt{978-0-8218-3515-4}. \textbf{DOI}: \texttt{10.1090/surv/110}.

\item[Perelman, G.] (2002). \textit{The Entropy Formula for the Ricci Flow and its Geometric Applications}. arXiv:math/0211159. \textbf{DOI}: \texttt{10.48550/arXiv.math/0211159}.
\end{description}

\vspace{12pt}
\begin{center}
\rule{0.5\textwidth}{0.4pt}
\end{center}
\vspace{6pt}

\textit{Nota bibliogr ifica:} Todos os DOIs foram verificados via API do CrossRef (\texttt{api.crossref.org}) e os ISBNs confirmados nas bases WorldCat (\texttt{worldcat.org}) e OpenLibrary (\texttt{openlibrary.org}). A data de verifica SS  o  c 30 de maio de 2026. Para \textit{preprints} hist 3ricos sem DOI (e.g., Ilinski \& Kalinin, 1998), indica-se o identificador arXiv. A URL de acesso aberto da monografia de Nelson (1967) foi inclu -da por cortesia, visto que o autor a disponibiliza livremente em seu \textit{site} institucional.

\endinput


